\documentclass[10pt,a4paper]{amsart}
\usepackage[margin=19mm]{geometry}
\usepackage{amsmath,amssymb,mathtools}
\usepackage[scr=rsfs,cal=euler]{mathalfa}
\usepackage{xfrac}
\usepackage[unicode,hidelinks,bookmarks=false]{hyperref}
\newcommand{\ie}{i.e.\ }
\newcommand{\cf}{cf.\ }
\newcommand{\resp}{respectively\ }
\newcommand{\Subsection}{\subsection}
\providecommand{\nobreakdash}{\nobreak\hbox{-}}
\setlength{\emergencystretch}{2em}
\multlinegap0pt
\usepackage{dsfont}
\usepackage{bm}
\usepackage{amssymb}
\usepackage{mathtools}
\usepackage{tikz}
\usepackage{xcolor}
\usepackage[shortlabels]{enumitem}
\newtheorem{theorem}{Theorem}
\newcommand{\PP}{\mathbb{P}}
\newcommand{\correction}[1]{{\textcolor{blue}{#1}}}
\newcommand{\e}{\mathrm{e}}
\title{Matrix Representations for Scale Functions of Spectrally Negative Lévy Processes with Rational Jumps}
\author{Osvaldo Angtuncio Hern\'andez, Oscar Peralta}
\date{Source: arXiv:2604.08889v1 (CC BY 4.0)}
\begin{document}
\maketitle

\noindent\emph{Source and licence.} Matrix Representations for Scale Functions of Spectrally Negative Lévy Processes with Rational Jumps. Version arXiv:2604.08889v1 (CC BY 4.0), distributed by arXiv under the Creative Commons Attribution 4.0 International licence, http://creativecommons.org/licenses/by/4.0/, as recorded in the arXiv licence metadata for this exact version and shown on the abstract page https://arxiv.org/abs/2604.08889v1. Full licence notice: \texttt{licenses/arxiv-2604.08889-CC-BY-4.0.txt}.\par\smallskip

\noindent\textit{Editorial scope.} Counted unit: the article's theorem teoScaleFnUV (source0.tex lines 265--292). The article's complete original proof of it is delivered with it (source0.tex lines 1103--1152, one \texttt{proof} environment of 50 lines, which the article places away from the statement). No mathematics is written, repaired or re-derived: the statement and the proof are the article's own lines, in the article's own notation, and no retained line is substituted or re-flowed.  Retained uncounted as the source-local context the proof needs: lines 178--180; lines 229--231; lines 548--550; lines 632--634; lines 636--638; lines 758--760; lines 882--889; lines 1089--1091; lines 1092--1094.  Nothing was omitted from any retained span, so the package carries no omission record.  No retained line contains a citation, so the wrapper carries no bibliography and no bibliography entry.  No retained line contains a figure environment, an image-inclusion command or drawing code, so no image is delivered and the package declares no asset dependency.  Editorial formatting only, in addition: an AMS \texttt{amsart} preamble that restates the article's own declarations and macros used by the retained lines, namely the environments \texttt{theorem} on separate counters, together with the standard packages those lines need.  The retained environments print in this document as Theorem 1 in order of appearance, whereas the article prints the same items with the numbering of its own full document.  The complete native source of the article as distributed in the arXiv e-print for this exact version is bundled unchanged as \texttt{sources/198177/source0.tex}.
\begin{equation}\label{eqnXTrajectories}
X_t=dt+\sigma B_t-\sum_{j=1}^{N_t}C_j,
\end{equation}with $\{C_j\}_j$  i.i.d. having distribution $F$, and $N$ is a Poisson process with rate $\lambda$.

\begin{equation}\label{eq:Ivanovs_identity}
W^{(q)}(x) = \frac{1}{\psi'(\Phi_q)} \Big( \e^{\Phi_q x} - \PP(\tau_{\{-x\}} < e_q) \Big), \qquad x \ge 0,
\end{equation}

\begin{equation}\label{eq:residual_sojourn_component}
\mathbb{P}(\text{residual sojourn time at } t > h \mid \bm{A}_t = (\bm{0}, \ldots, \bm{0}, \bm{a}, \bm{0}, \ldots, \bm{0})) = \bm{a} e^{\bm{C}_{ii}^k h} \bm{1} \ge 0, \qquad h \ge 0.
\end{equation}

\begin{equation}\label{eq:O_x_formula_minus}
\mathbb{E}_{\bm{\beta}}[\bm{O}_x \mathds{1}\{\gamma_{\{-x\}} < \infty\}] = \bm{\beta} \, e^{(\bm{C}^- + \bm{D}^{-+}\bm{\Psi})x}.
\end{equation}

\begin{equation}\label{eq:O_x_formula}
\mathbb{E}_{\bm{\alpha}}[\bm{O}_x \mathds{1}\{\gamma_{\{-x\}} < \infty\}] = \bm{\alpha} \bm{\Psi} \, e^{(\bm{C}^- + \bm{D}^{-+}\bm{\Psi})x}.
\end{equation}

\begin{equation}\label{eq:hitting_factor}
\PP(\tau_{\{-x\}} < e_q) = \mathbb{E}\Big[\mathds{1}\big\{\tau^-_{\{-x\}} < e_q\big\}\mathbb{E}_{-Z}\big[\mathds{1}\{\tau_{\{0\}} < e_q\}\big]\Big] = \mathbb{E}\Big[\e^{-\Phi_q Z}\;\mathds{1}\big\{\tau^-_{\{-x\}} < e_q\big\}\Big].
\end{equation}

\begin{equation}\label{eqnBoundedVariationTauxProb}
\begin{split}
\PP(\tau_{\{-x\}} < e_q) 
&= \mathbb{E}\Big[\e^{-\Phi_q Z}\;\mathds{1}\{\gamma_{\{-x\}} < \infty\}\Big]\\
&= \mathbb{E}_{\bm A_0}\Big[\bm{O}_x \bm{\nu}\;\mathds{1}\{\gamma_{\{-x\}} < \infty\}\Big]\\
&= \mathbb{E}_{\bm A_0}\Big[\bm{O}_x\;\mathds{1}\{\gamma_{\{-x\}} < \infty\}\Big]\bm{\nu}.
\end{split}
\end{equation}Starting from $\bm{A}_0 = 1 \in \mathcal{Z}^+$, the downward record formula \eqref{eq:O_x_formula} gives $\mathbb{E}_{\bm A_0}[\bm{O}_x\;\mathds{1}\{\gamma_{\{-x\}} < \infty\}] = \bm{\Psi}\, e^{\bm{G}x}$, where $\bm{G} := \bm{C}^- + \bm{D}^{-+}\bm{\Psi}= \bm{T} + \bm{t}\,\bm{\Psi}$. Therefore,

\begin{equation}\label{eq:params_ubv}
\bm{C}^- = \begin{pmatrix} -\omega_{\lambda+q} & \bm{0} \\ \bm{t} & \bm{T} \end{pmatrix}, \qquad \bm{C}^+ = -\eta_{\lambda+q},
\end{equation}

\begin{equation}\label{eq:params_ubv1}
\bm{D}^{-+} = \begin{pmatrix} \omega_{\lambda+q} \\ \bm{0} \end{pmatrix}, \qquad \bm{D}^{+-} = \begin{pmatrix} 0 & \eta_{\lambda+q}\dfrac{\lambda}{\lambda+q}\,\bm{\alpha} \end{pmatrix},
\end{equation}

\begin{theorem}\label{teoScaleFnUV}
Let $X$ be a L\'evy process as in \eqref{eqnXTrajectories}, with $\sigma>0$ and $d\in \mathbb{R}$. 
Consider $\bm{G}: = \bm{C}^- + \bm{D}^{-+}\bm{\Psi}$ which, under the block structure given below in \eqref{eq:params_ubv}-\eqref{eq:params_ubv1}, takes the explicit form $\bm{G} = \begin{pmatrix} -a & \bm{b} \\ \bm{t} & \bm{T} \end{pmatrix}$. Define $\bm{\nu}:=(\Phi_q \bm{I}-\bm{T})^{-1}\bm{t}$ and $\bm{V}:=\binom{1}{\bm{\nu}}$. Let $\bm{e}_1^-:=(1,0,\dots,0)^\intercal\in\mathbb{R}^{p+1}$ be the first canonical basis vector.
Then, for any $x\geq 0$ and $q\geq 0$ such that $\psi'(\Phi_q)>0$ we have
\begin{equation}\label{eq:scale_ubv}
W^{(q)}(x) = \frac{1}{\psi'(\Phi_q)}\Big(e^{\Phi_q x} - (\bm{e}_1^-)^\intercal\,e^{\bm{G}x}\,\bm{V}\Big), \qquad x \ge 0.
\end{equation}
The matrix $\bm{\Psi}$ can be 
recovered from the pair $(a,\bm{b})$ via
\begin{equation}\label{eq:Psi_from_ab}
\bm{\Psi} = \frac{1}{\omega_{\lambda+q}}\begin{pmatrix}\omega_{\lambda+q} - a & \quad \bm{b}\end{pmatrix},
\end{equation}
where $\omega_{\lambda+q} := \big(\sqrt{d^2 + 2\sigma^2 (\lambda+q)} + d\big)\sigma^{-2}$, $\eta_{\lambda+q} := \big(\sqrt{d^2 + 2\sigma^2 (\lambda+q)} - d\big)\sigma^{-2}$, and $(a,\bm b)$ are a solution to the system of equations
\begin{equation}\label{eq:ivanovs_scalar}
\sigma^2a^2-2da-2(\lambda+q)+\sigma^2\bm b\bm t=0,
\end{equation}
\begin{equation}\label{eq:ivanovs_vector}
\sigma^2\bm{b}\left(\big(a-2d\sigma^{-2}\big)\bm I-\bm{T}\right) =2\lambda\bm{\alpha}.
\end{equation}
Finally, such a pair can be obtained as the limit of the recursive equations
\begin{equation}\label{eq:b_recursion}
\bm{b}_n = \Big( \frac{2\lambda}{\sigma^2}\,\bm{\alpha} + (\omega_{\lambda+q}-a_{n-1})\bm{b}_{n-1} \Big) \big(\eta_{\lambda+q}\bm{I} - \bm{T}\big)^{-1}, \qquad n \ge 1,
\end{equation}
\begin{equation}\label{eq:a_recursion}
a_n = \omega_{\lambda+q} - \frac{\sigma^2}{2\sqrt{d^2+2\sigma^2(\lambda+q)}}\Big((\omega_{\lambda+q} - a_{n-1})^2 + \bm{b}_n\bm{t}\Big), \qquad n \ge 1,
\end{equation}
with $a_0=\omega_{\lambda+q}$, $\bm{b}_0=\bm{0}$.
\end{theorem}

\begin{proof}[Proof of first part of Theorem \ref{teoScaleFnUV}]
Let $\bm{e}_1^-$ denote the unit column vector of dimension $p+1$ with one in the first entry, so that $(\bm{e}_1^-)^\intercal$ selects the initial Brownian downward state. Under the block ordering $(\mathcal{Z}_1^-,\mathcal{Z}_2^-)$, the region $\mathcal{Z}_1^-$ is represented by the singleton state $\{(\bm{e}_1^-)^\intercal\}$, whereas $\mathcal{Z}_2^-$ consists of a subset of all row vectors of the form $(0,\bm{a})$ with $\bm{a}\in\mathbb{R}^{1\times p}$ and $\bm{a}\bm{1}=1$. Thus the first coordinate corresponds to the Brownian downward phase, and the remaining $p$ coordinates correspond to the matrix-exponential drain phase.
Note that $\{\tau^-_{\{-x\}} < e_q\} = \{\gamma_{\{-x\}} < \infty\}$ in the RAP-modulated embedding, since the $q$-killing is encoded as termination from $\mathcal{Z}^+$: the orbit reaches level $-x$ if and only if the stage when it terminates happens after reaching $-x$, and the latter happens if and only if the excursion where the L\'evy process is killed occurs after hitting $(-\infty,-x)$. 
%Note that the $q$-killing of the L\'evy process occurs between times $(\tau_n,\tau_{n+1})$ for some $n$, if and only if the RAP terminates at the end of the $n$-th excursion. }
%\correction{I added this because we are not being explicit on the correspondence between the excursion with the $q$-killing and the RAP terminating. The reader might ask herself what happens if in the same excursion the Levy hits -x AND dies? IT is not a problem since at the END of the $n$-th excursion both die}. 
%The dependence on $q$ is absorbed into the rates $\omega_{\lambda+q}$ and $\eta_{\lambda+q}$ via the Wiener-Hopf factorization. 
This equality is realized under the stagewise coupling induced by the Wiener-Hopf decomposition. 
%The fluid stage and the corresponding Brownian stage start at the same level, have the same running minimum, end at the same pre-jump level, and carry the same jump-versus-killing mark. Consequently, a given stage downcrosses $-x$ before termination in the fluid picture if and only if the corresponding Brownian stage downcrosses $-x$ before the killing time in the Lévy picture. Iterating this comparison stage by stage yields the event identity $\{\tau^-_{\{-x\}}<e_q\}=\{\gamma_{\{-x\}}<\infty\}$. 


We use \eqref{eq:hitting_factor} and decompose on the events $\{Z=0\}$ and $\{Z>0\}$, obtaining
\[
\PP(\tau_{\{-x\}} < e_q) = \mathbb{P}\big(\tau^-_{\{-x\}} < e_q,Z=0\big)+\mathbb{E}\Big[\e^{-\Phi_q Z}\;\mathds{1}\big\{\tau^-_{\{-x\}} < e_q,Z>0\big\}\Big].
\]
When $Z=0$ the process hits $-x$ by a downward Brownian excursion, so $\mathds{1}\{Z=0\}=\bm O_{x}\bm{e}_1^-$, and thus
\[
\mathbb{P}\big(\tau^-_{\{-x\}} < e_q,Z=0\big)=\mathbb{E}_{(\bm{e}_1^-)^\intercal}\big[\bm O_x\mathds{1}\{\gamma_{\{-x\}}<\infty\}\big]\bm{e}_1^-.
\]
From \eqref{eq:O_x_formula_minus}, we obtain
\[
\mathbb{P}\big(\tau^-_{\{-x\}} < e_q,Z=0\big)=(\bm{e}_1^-)^\intercal e^{\bm{G}x}\bm{e}_1^-,
\]
where $\bm{G}=\bm{C}^- + \bm{D}^{-+}\bm{\Psi}$.

On the other hand, $\{\tau^-_{\{-x\}} < e_q, Z>0\} = \{\gamma_{\{-x\}}<\infty, \bm{O}_x\in \mathcal{Z}^-_2\}$. Given $\bm{O}_x = (0,\bm{a}) \in \mathcal{Z}_2^-$ with $\bm{a} \in \mathbb{R}^{1\times p}$, the overshoot $Z$ is the residual sojourn time in $\mathcal{Z}_2^-$. Since $\bm{C}_{22}^- = \bm{T}$, Equation \eqref{eq:residual_sojourn_component} gives the survival function $\bm{a}e^{\bm{T}z}\bm{1}$. By the same computation as in the bounded variation proof \eqref{eqnBoundedVariationTauxProb},
\[
\mathbb{E}_{(\bm{e}_1^-)^\intercal}\big[\e^{-\Phi_q Z} \mid \bm{O}_x = (0,\bm{a})\big] = \bm{a}(\Phi_q\bm{I}-\bm{T})^{-1}\bm{t} = \bm{a}\,\bm{\nu}.
\]
To express this in terms of the full $(p+1)$-dimensional orbit vector $\bm{O}_x = (0,\bm{a})$, we use the identity matrix and the block extraction
\[
\bm{a} = (0,\bm{a})\begin{pmatrix}\bm{0}\\\bm{I}_p\end{pmatrix},
\]
so that
\[
\mathbb{E}_{(\bm{e}_1^-)^\intercal}\big[\e^{-\Phi_q Z} \mid \bm{O}_x=(0,\bm{a})\big] = (0,\bm{a})\begin{pmatrix}\bm{0}\\\bm{I}_p\end{pmatrix}\bm{\nu} = (0,\bm{a})\begin{pmatrix}0\\\bm{\nu}\end{pmatrix}.
\]
Thus,
\[
\mathbb{E}\Big[\e^{-\Phi_q Z}\;\mathds{1}\big\{\tau^-_{\{-x\}} < e_q,Z>0\big\}\Big] = \mathbb{E}_{(\bm{e}_1^-)^\intercal}\Big[\bm{O}_x\begin{pmatrix}0\\\bm{\nu}\end{pmatrix}\mathds{1}\{\gamma_{\{-x\}} < \infty, \bm{O}_x\in\mathcal{Z}_2^-\}\Big].
\]
Since on the event $\{\bm{O}_x\in\mathcal{Z}_2^-\}$ we have $\bm{O}_x=(0,\bm{a})$ with $\bm{a}\bm{1}=1$, and on the complementary event $\{\bm{O}_x\in\mathcal{Z}_1^-\}$ we have $\bm{O}_x\begin{pmatrix}0\\\bm{\nu}\end{pmatrix}=0$, the indicator $\mathds{1}\{\bm{O}_x\in\mathcal{Z}_2^-\}$ can be dropped and we may write
\[
\mathbb{E}\Big[\e^{-\Phi_q Z}\;\mathds{1}\big\{\tau^-_{\{-x\}} < e_q,Z>0\big\}\Big] = \mathbb{E}_{(\bm{e}_1^-)^\intercal}\Big[\bm{O}_x\,\mathds{1}\{\gamma_{\{-x\}} < \infty\}\Big]\begin{pmatrix}0\\\bm{\nu}\end{pmatrix} = (\bm{e}_1^-)^\intercal e^{\bm{G}x}\begin{pmatrix}0\\\bm{\nu}\end{pmatrix}.
\]
Combining both cases and setting $\bm{V}:=\begin{pmatrix}1\\\bm{\nu}\end{pmatrix}\in\mathbb{R}^{p+1}$, we obtain
\begin{equation}\label{eq:hitting_ubv}
\PP(\tau_{\{-x\}} < e_q) = (\bm{e}_1^-)^\intercal e^{\bm{G}x}\bm{V}, \qquad x \ge 0.
\end{equation}
The latter and \eqref{eq:Ivanovs_identity} imply \eqref{eq:scale_ubv}.
\end{proof}

\end{document}
